% from ega4-9-fr.tex
\oldpage[IV-3]{65}

cohérents $\mathcal G$, $\mathcal H$ (où $Z=W\cap h^{-1}(U_0)$,
$h:Y\to S$ étant le morphisme structural) et un $\mathcal O_Z$-homomorphisme
$u:\mathcal G\to\mathcal H$ tel que $\mathcal M=\mathcal G_\eta$,
$\mathcal N=\mathcal H_\eta$ et $v=u_\eta$~; on peut donc supposer $U_0$ pris
tel que, pour $s\in U_0$, $u_s:\mathcal G_s\to\mathcal H_s$ soit injectif
(9.4.5). Mais pour tout $s\in U_0$, l'homomorphisme
$1\otimes u_s:\mathcal F_s\otimes_{\mathcal O_{Z_s}}\mathcal G_s\to
\mathcal F_s\otimes_{\mathcal O_{Z_s}}\mathcal H_s$ n'est autre que
$(1\otimes u)_s$~; l'hypothèse que $(1\otimes u)_\eta$ est non injectif
entraîne donc (9.4.5) l'existence d'un ouvert non vide $U\subset U_0$ tel que
pour tout $s\in U$, $(1\otimes u)_s$ soit non injectif, et par suite
$\mathcal F_s$ n'est pas $g_s$-plat pour tout $s\in U$.

$4^\circ$ $\eta\in S-E'$. Il est clair que $S-E\subset S-E'$, et si
$\eta\in S-E$, $S-E'$ est \emph{a fortiori} un voisinage de $\eta$ d'après le
$3^\circ$. Supposons donc $\eta\in E$, donc $\mathcal F_\eta$ localement
libre~; ces hypothèses entraînent que $X_\eta$ est non connexe, et que les
rangs du $\mathcal O_{X_\eta}$-Module localement libre $\mathcal F_\eta$ ne
sont pas les mêmes sur les diverses composantes connexes de $X_\eta$. Or, il
résulte de (8.4.2), appliqué suivant la méthode de (8.1.2, \emph{a)}), que l'on
peut supposer (en remplaçant $S$ par un voisinage ouvert de $\eta$) que $X$ et
$X_\eta$ ont le même nombre de composantes connexes, les composantes connexes
de $X_\eta$ étant les intersections de $X_\eta$ et des composantes connexes de
$X$. La conclusion résulte alors du raisonnement fait dans le $2^\circ$,
appliqué à chacune des composantes connexes de $X$ (qui sont en nombre fini).

\begin{remark}[9.4.7.2]
\label{IV.9.4.7.2-fr}
Le lemme (9.4.7.1) sera plus tard généralisé et débarrassé d'hypothèses
noethériennes (11.2.8).
\end{remark}

\begin{proposition}[9.4.8]
\label{IV.9.4.8-fr}
Soient $S$ un préschéma localement noethérien, $f:X\to S$ un morphisme de type
fini, $\mathcal F$ un $\mathcal O_X$-Module cohérent. On suppose que pour tout
$s\in S$, $X_s$ soit un préschéma localement intègre. Alors l'ensemble $E$
\textnormal{(}resp. $E'$\textnormal{)} des $s\in S$ tels que $\mathcal F_s$
soit un $\mathcal O_{X_s}$-Module de torsion \textnormal{(}resp. un
$\mathcal O_{X_s}$-Module sans torsion\textnormal{)} est localement
constructible.
\end{proposition}

On peut évidemment supposer $S=\operatorname{Spec}(A)$ affine et noethérien et
prouver que $E$ (resp. $E'$) est alors constructible en utilisant le critère
$(\mathbf 0_{\mathrm{III}},9.2.3)$~; remplaçant $S$ par le sous-préschéma fermé
réduit de $S$ ayant pour espace sous-jacent une partie fermée irréductible $Y$
de $S$, et notant (\textbf{I}, 3.6.4) que pour $s\in Y$, la fibre
$(X_{(Y)})_s$ s'identifie canoniquement à $X_s$ et le faisceau
$(\mathcal F\otimes_{\mathcal O_S}\mathcal O_Y)_s$ à $\mathcal F_s$, on voit
qu'on est ramené au cas où $S$ est intègre de point générique $\eta$, et à
prouver que $E$ ou $S-E$ (resp. $E'$ ou $S-E'$) est un voisinage de $\eta$
dans $S$. Notons en outre que $X$ est réunion finie d'ouverts affines $U_i$,
de type fini sur $S$, et chacun des $(U_i)_s$ est induit sur un ouvert de
$X_s$, donc localement intègre~; en outre, si $(U_i)_\eta$ est vide (resp. non
vide), on sait que $(U_i)_s$ est aussi vide (resp. non vide) dans un voisinage
de $\eta$ (9.2.6). On peut donc supposer tous les $(U_i)_\eta$ non vides et
intègres, et dire que $\mathcal F_s$ est de torsion (resp. sans torsion)
équivaut à dire que chacun des $\mathcal F_s|(U_i)_s$ est de torsion (resp.
sans torsion). On est donc ramené au cas où $X=\operatorname{Spec}(B)$ est
affine, et $\mathcal F=\widetilde M$, où $M$ est un $B$-module de type fini~;
on posera $B_s=B\otimes_A k(s)$, $M_s=M\otimes_A k(s)$ et l'on peut supposer
$B_\eta$ intègre. Nous avons quatre cas à envisager~:

$1^\circ$ $\eta\in E$~; $M_\eta$ est donc un $B_\eta$-module de torsion de
type fini, et il y a par suite un $h\neq0$ dans $B_\eta$ tel que
$hM_\eta=0$~; en vertu de (8.5.2, (i)), appliqué suivant la méthode

